\documentclass[12pt, a4paper]{article}
\usepackage[utf8]{inputenc}
\usepackage[T2A]{fontenc}
\usepackage{indentfirst, setspace}
\usepackage{tabularx, multirow}
\usepackage[normalem]{ulem}
\usepackage[style=russian]{csquotes}
\usepackage[english,russian]{babel}
\usepackage{hyperref}
\usepackage{ragged2e}
\usepackage{caption}
\usepackage{wrapfig}
\usepackage{amsmath}
\usepackage{tikz}
\makeatletter
\def\@biblabel#1{#1. }
\makeatother
\captionsetup{labelsep=endash}
\usepackage{listings}
\linespread{1.3}
\lstset{
  language=C++, % Язык программирования (можно выбрать другой)
  basicstyle=\ttfamily\small, % Основной стиль шрифта
  keywordstyle=\color{blue}, % Стиль ключевых слов
  breaklines=true,
  commentstyle=\color{green}, % Стиль комментариев
  stringstyle=\color{red}, % Стиль строк
  extendedchars=\true,
  showstringspaces=false,
}
\usepackage[left=20mm, top=15mm, right=15mm, bottom=20mm]{geometry}

\begin{document}
\begin{titlepage}
     \fontsize{12}{12}\selectfont

  {\centering

   \begin{bf}

    \begin{wrapfigure}{l}{19mm}
        
\includegraphics[width=22mm]{emblem.jpeg}
    \end{wrapfigure}


    \noindent Министерство науки и высшего образования Российской Федерации

    \noindent Федеральное государственное бюджетное образовательное учреждение высшего образования

    \noindent \enquote{Московский государственный технический унверситет

     \noindent имени Н.Э. Баумана

     \noindent (национальный исследовательский университет)}

    \noindent (МГТУ им. Н.Э. Баумана)

   \end{bf}
  }

  \vspace{0.4cm}

  {\setstretch{0.1}
   \noindent\rule{\textwidth}{1mm}
   \noindent\rule{\textwidth}{0.5mm}

  }

  \fontsize{14}{21}\selectfont

  \noindent\begin{tabularx}{\textwidth}{l >{\centering\arraybackslash}X}
   ФАКУЛЬТЕТ & \flqq Фундаментальные Науки\frqq \\ \cline{2-2}

   КАФЕДРА & ФН-12 \flqq Математическое моделирование\frqq \\ \cline{2-2}
  \end{tabularx}


  \vspace{1cm}


  \begin{center}
   \begin{bf}

    \fontsize{24}{36}\selectfont
    ОТЧЕТ

    \fontsize{20}{30}\selectfont
    ПО ЛАБОРАТОРНОЙ РАБОТЕ НА ТЕМУ:

    Длинная арифметика

   \end{bf}
  \end{center}

  \fontsize{14}{21}\selectfont
  \vspace{5cm}


  \noindent\begin{tabularx}{\textwidth}{ X >{\centering}p{4cm} p{1cm} c }
   Студент: & & & Подобедов В. В. \\ \cline{2-2} \cline{4-4}
   & \fontsize{10}{15}\selectfont дата, подпись & & \fontsize{10}{15}\selectfont Ф.И.О. \\
   Преподаватель: & & & Волкова Л. Л.\\ \cline{2-2} \cline{4-4}
   & \fontsize{10}{15}\selectfont дата, подпись & & \fontsize{10}{15}\selectfont Ф.И.О.
  \end{tabularx}

  \vspace{\fill}

  \begin{center}
   \it{Москва}, 2023
  \end{center}

  \thispagestyle{empty}
\end{titlepage}

\renewcommand{\contentsname}{Содержание}
\setcounter{page}{2}
\tableofcontents

\newpage
\begin{center}
   \section*{Введение} 
\addcontentsline{toc}{section}{Введение}
\end{center}
\justifying

\textbf{Цель лабораторной работы}: создать собственную структуру данных.

Для достижения поставленной цели требуется решить следующие \textbf{задачи}.
\begin{enumerate}
\item Описать троичные числа.
\item Разработать структуру данных для хранения целых чисел с повышением ёмкости хранения по сравнению с базовым типом данных.
\item Реализовать $6$ операторов сравнения, операторы присваивания (копии значения), сложения и вычитания двух структур длинных чисел.
\item Выполнить тестирование.
\end{enumerate}

Согласно варианту, требуется разработать структуру хранения троичных чисел.

\newpage
\begin{center}
   \section{Аналитическая часть} 
\end{center}

\textbf{Длинная арифметика} --- арифметика, которая относится к арифметическим операциям над числами, превышающими предельные значения стандартных числовых типов данных (например, ``int'', ``long'', ``double'' и т.д.), либо к работе с числами, которые требуют более высокой точности или специфических условий хранения.

\textbf{Троичные числа}~\cite{book} --- это числа, которые используют систему счисления с основанием $3$. В отличие от привычной для нас десятичной системы счисления (основание $10$), где используются цифры от $0$ до $9$, в троичной системе используются только три цифры: $0$, $1$ и $2$.

\textbf{Емкость хранения} --- количество данных или информации, которое может быть сохранено или содержано в определенной системе, устройстве, структуре данных или носителе информации.

\textbf{Тип данных} --- определяет характеристики данных, их формат, допустимый диапазон значений и операции, которые можно выполнять над этими данными.

\newpage
\begin{center}
    \section{Конструкторская часть}
\end{center}
\subsection{Написание собственной структуры данных}
С помощью класса ``LongTritNumber'', который описывает структуру данных для хранения целых чисел с повышенной емкостью хранения по сравнению с базовым типом данных, реализовать арифметику и сравнения для чисел, представленных троичными цифрами (False, Unknown, True) с емкостью хранения в 8 тритов. Требуется реализовать следующие методы и перегруженные операторы.
\begin{itemize}
    \item[---] Метод ``addTrits'' --- метод, который складывает два трита и возвращает сумму.
    \item[---] Метод ``calculateCarry'' --- метод, который вычисляет перенос при сложении троичных чисел.
    \item[---] Метод ``subtractTrits'' --- метод, который вычитает два трита и возвращает разность.
    \item[---] Метод ``calculateBorrow'' --- метод, который вычисляет перенос при вычитании троичных чисел.
    \item[---] ``operator[]'' --- перегруженный оператор для доступа к тритам по индексу.
    \item[---] ``operator='' --- перегруженный оператор для операции присваивания.
    \item[---] Операторы сравнения (``operator=='', ``operator!='', ``operator<'', ``operator>'', ``operator<='', ``operator>='') --- реализуют сравнение двух чисел по их тритам.
    \item[---] Операторы арифметики (``operator+'', ``operator-'') --- реализуют сложение и вычитание двух чисел.
    \item[---] Дружественная функция-оператор ``operator<<'' --- перегруженный оператор для вывода числа в виде строки троичных значений.
    \item[---] Дружественная функция-оператор ``operator>>'' --- перегруженный оператор для ввода числа с клавиатуры.
\end{itemize}

\subsection{Сложение двух чисел}
\begin{enumerate}
    \item Создается новый объект типа ``LongTritNumber'', предназначенный для хранения результата сложения.
    \item Инициализируется переменная ``carry'' с начальным значением False.
    \item Итерация по тритам.
    \begin{itemize}
        \item[---] Происходит итерация через триты чисел с конца к началу.
        \item[---] Выполняется вызов вспомогательной функции ``addTrits'' для сложения тритов и получения суммы и значения переноса.
        \item[---] Результат сложения помещается в соответствующий разряд нового объекта типа ``LongTritNumber''.
        \item[---] Перенос сохраняется для использования на следующей итерации.
    \end{itemize}
    \item  Возвращается объект типа ``LongTritNumber'', содержащий результат сложения.
\end{enumerate}

\subsection{Вычитание двух чисел}
\begin{enumerate}
    \item Создается новый объект типа ``LongTritNumber'', предназначенный для хранения результата вычитания.
    \item Инициализируется переменная ``borrow'' с начальным значением False.
    \item Итерация по тритам.
    \begin{itemize}
        \item[---] Происходит итерация через триты чисел с конца к началу.
        \item[---] Выполняется вызов вспомогательной функции ``subtractTrits'' для вычитания тритов и получения разности и значения переноса.
        \item[---] Результат вычитания помещается в соответствующий разряд нового объекта типа ``LongTritNumber''.
        \item[---] Перенос сохраняется для использования на следующей итерации.
    \end{itemize}
    \item  Возвращается объект типа ``LongTritNumber'', содержащий результат вычитания.
\end{enumerate}

\subsection{Операторы сравения}
\begin{enumerate}
    \item Оператор равенства.
    \begin{itemize}
        \item[---] Пройти по всем тритам в обоих числах, сравнивая их между собой по позициям.
        \item[---] Если все триты в соответствующих позициях равны, возвращает ``true''. Если хотя бы один трит не совпадает, возвращает ``false''.
    \end{itemize}
    \item Оператор неравенства.
    \begin{itemize}
        \item[---] Пройти по всем тритам в обоих числах, сравнивая их между собой по позициям.
        \item[---] Если все триты в соответствующих позициях равны, возвращает ``false''. Если хотя бы один трит не совпадает, возвращает ``true''.
    \end{itemize}
    \item Оператор сравнения по отношению.
    \begin{itemize}
        \item[---] Пройти по тритам чисел, начиная с наиболее значимых разрядов
        \item[---] Если найдено отличие, операторы возвращают ``true'' или ``false'' в соответствии с результатом сравнения.
        \item[---] Если все триты в соответствующих позициях равны, операторы могут продолжить сравнение следующих тритов, пока не будет найдено отличие или не закончатся триты для сравнения. 
    \end{itemize}
\end{enumerate}

\subsection{Оператор ввода}
\begin{enumerate}
    \item Для каждого элемента в векторе ``trits'' (переменная, хранящая триты объекта\\ ``LongTritNumber'') выполняется итерация. При каждой итерации запрашивается у пользователя значение для трита с помощью стандартного ввода.
    \item После ввода значения проверяется, находится ли оно в допустимом диапазоне ($0$, $1$ или $2$). Это происходит с помощью проверки условия ``if''. Если введенное значение не соответствует триту ($0$, $1$ или $2)$, выводится сообщение о некорректном вводе и оставшиеся незаполненны элементы вектора автоматически заполняются значениями $0$. Если значение корректно, оно сохраняется в текущий трит в векторе ``trits''.
    \item После завершения итерации по всем тритам в объекте ``LongTritNumber'' функция возвращает поток ввода.
\end{enumerate}

\subsection{Оператор вывода}
\begin{enumerate}
    \item Для каждого элемента в векторе ``trits'' (переменная, хранящая триты объекта\\ ``LongTritNumber'') выполняется итерация. При каждой итерации определяется значение трита и выводится соответствующее ему текстовое представление в поток вывода.
    \item После завершения итерации по всем тритам в объекте ``LongTritNumber'' функция возвращает поток вывода.
\end{enumerate}

\newpage
\begin{center}
    \section{Технологическая часть}
\end{center}

Лабораторная работа была реализована на языке программирования C++. Кроме стандартной библиотеки <iostream>, которая предоставляет возможность пользоваься средствами ввода-вывода для стандартной консоли, была подключена библиотека <vector>. Она представляет собой динамический массив, способный изменять свой размер в процессе выполнения программы.

Далее приведен листинг реализации стека (листинг~\ref{lst:label1}).
\begin{lstlisting}[caption={Исходный код программы для реализации стека},label={lst:label1}]
//#include <iostream>
//#include <vector>
//
//using namespace std;
//
//class LongTritNumber {
//public:
//	enum class Trit {False = 0, Unknown = 1, True = 2 };
//	static const int Capacity = 8; 
//
//	LongTritNumber() {
//		trits.resize(Capacity, Trit::False); 
//	}
//
//	const Trit& operator[](size_t index) const {
//		return trits[index];
//	}
//
//	LongTritNumber& operator=(const LongTritNumber& other) {
//		if (this != &other) {
//			for (size_t i = 0; i < Capacity; ++i) {
//				trits[i] = other.trits[i];
//			}
//		}
//		return *this;
//	}
//
//	bool operator==(const LongTritNumber& other) const {
//		for (size_t i = 0; i < Capacity; ++i) {
//			if (trits[i] != other.trits[i]) {
//				return false;
//			}
//		}
//		return true;
//	}
//
//	bool operator!=(const LongTritNumber& other) const {
//		for (size_t i = 0; i < Capacity; ++i) {
//			if (trits[i] == other.trits[i]) {
//				return true;
//			}
//		}
//		return true;
//	}
//
//	LongTritNumber operator+(const LongTritNumber& other) const {
//		LongTritNumber result;
//		Trit carry = Trit::False;
//
//		for (int i = Capacity - 1; i >= 0; --i) {
//			Trit sum = addTrits(trits[i], other.trits[i], carry);
//			result.trits[i] = sum;
//			carry = calculateCarry(trits[i], other.trits[i], carry);
//		}
//
//		return result;
//	}
//
//	LongTritNumber operator-(const LongTritNumber& other) const {
//		LongTritNumber result;
//		Trit borrow = Trit::False;
//
//		for (int i = Capacity - 1; i >= 0; --i) {
//			Trit diff = subtractTrits(trits[i], other.trits[i], borrow);
//			result.trits[i] = diff;
//			borrow = calculateBorrow(trits[i], other.trits[i], borrow);
//		}
//
//		return result;
//	}
//
//	bool operator<(const LongTritNumber& other) const {
//		for (size_t i = 0; i < Capacity; ++i) {
//			if (trits[i] < other.trits[i]) {
//				return true;
//			}
//			else if (trits[i] > other.trits[i]) {
//				return false;
//			}
//		}
//		return false;
//	}
//
//	bool operator>(const LongTritNumber& other) const {
//		for (size_t i = 0; i < Capacity; ++i) {
//			if (trits[i] > other.trits[i]) {
//				return true;
//			}
//			else if (trits[i] < other.trits[i]) {
//				return false;
//			}
//		}
//		return false;
//	}
//
//	bool operator<=(const LongTritNumber& other) const {
//		for (size_t i = 0; i < Capacity; ++i) {
//			if (trits[i] <= other.trits[i]) {
//				return true;
//			}
//			else if (trits[i] >= other.trits[i]) {
//				return false;
//			}
//		}
//		return false;
//	}
//
//	bool operator>=(const LongTritNumber& other) const {
//		for (size_t i = 0; i < Capacity; ++i) {
//			if (trits[i] >= other.trits[i]) {
//				return true;
//			}
//			else if (trits[i] <= other.trits[i]) {
//				return false;
//			}
//		}
//		return false;
//	}
//
//	friend ostream& operator<<(ostream& os, const LongTritNumber& number) {
//		for (const auto& trit : number.trits) {
//			switch (trit) {
//			case Trit::False:
//				os << "False";
//				break;
//			case Trit::Unknown:
//				os << "Unknown";
//				break;
//			case Trit::True:
//				os << "True";
//				break;
//			}
//			os << " ";
//		}
//		return os;
//	}
//
//	friend istream& operator>>(istream& is, LongTritNumber& number) {
//		for (auto& trit : number.trits) {
//			int userInput;
//			cout << "Введите значение трита (0 - False, 1 - Unknown, 2 - True): ";
//			is >> userInput;
//
//			if (userInput >= 0 && userInput <= 2) {
//				trit = static_cast<Trit>(userInput);
//			}
//			else {
//				cout << "Некорректный ввод, используйте значения от 0 до 2." << endl;
//				return is;
//			}
//		}
//		return is;
//	}
//
//private:
//	vector<Trit> trits;
//
//	Trit addTrits(Trit a, Trit b, Trit& carry) const {
//		int sum = static_cast<int>(a) + static_cast<int>(b);
//		carry = static_cast<Trit>(sum / 3);
//		return static_cast<Trit>(sum % 3);
//	}
//
//	Trit calculateCarry(Trit a, Trit b, Trit carry) const {
//		int sum = static_cast<int>(a) + static_cast<int>(b) + static_cast<int>(carry);
//		return static_cast<Trit>(sum / 3); // 1 если sum >= 3, иначе 0
//	}
//
//	Trit subtractTrits(Trit a, Trit b, Trit& borrow) const {
//		int diff = static_cast<int>(a) - static_cast<int>(b) - static_cast<int>(borrow);
//		if (diff < 0) {
//			diff += 3; // Коррекция для отрицательных значений
//			borrow = Trit::True;
//		}
//		else {
//			borrow = Trit::False;
//		}
//		return static_cast<Trit>(diff);
//	}
//
//	Trit calculateBorrow(Trit a, Trit b, Trit borrow) const {
//		int diff = static_cast<int>(a) - static_cast<int>(b) - static_cast<int>(borrow);
//		return (diff < 0) ? Trit::True : Trit::False; // 1 если diff < 0, иначе 0
//	}
//};
//
//int main() {
//	setlocale(0, "");
//	LongTritNumber a, b; 
//
//	int otr1 = 0;
//	cout << "Число а положительное или отрицательное? (0 - полож. / 1 - отриц)";
//	cin >> otr1;
//    cout << "Введите значения для переменной a:" << endl;
//    cin >> a;
//
//	int otr2 = 0;
//	cout << "Число b положительное или отрицательное? (0 - полож. / 1 - отриц)";
//	cin >> otr2;
//    cout << "Введите значения для переменной b:" << endl;
//    cin >> b;
//
//	LongTritNumber c = a + b; 
//	LongTritNumber d = a - b; 
//
//	cout << "Результат сложения: " << c << endl;
//	cout << "Результат вычитания: " << d << endl;
//
//	if (a == b) {
//		cout << "a равно b" << endl;
//	}
//	else {
//		cout << "a не равно b" << endl;
//	}
//
//	if (a != b) {
//		cout << "a не равно b" << endl;
//	}
//	else {
//		cout << "a равно b" << endl;
//	}
//
//	if (a < b) {
//		cout << "a меньше b" << endl;
//	}
//	else {
//		cout << "a не меньше b" << endl;
//	}
//
//	if (a > b) {
//		cout << "a больше b" << endl;
//	}
//	else {
//		cout << "a не больше b" << endl;
//	}
//
//	if (a <= b) {
//		cout << "a меньше или равно b" << endl;
//	}
//	else {
//		cout << "a не меньше или равно b" << endl;
//	}
//
//	if (a >= b) {
//		cout << "a больше или равно b" << endl;
//	}
//	else {
//		cout << "a не больше или равно b" << endl;
//	}
//
//	return 0;
//}
//


#include <iostream>
#include <vector>

using namespace std;

class LongTritNumber {
public:
	enum class Trit { False = 0, Unknown = 1, True = 2 };
	static const int Capacity = 9; // Увеличиваем емкость на 1 для хранения знака

	LongTritNumber() {
		trits.resize(Capacity, Trit::False);
	}

	const Trit& operator[](size_t index) const {
		return trits[index];
	}

	LongTritNumber& operator=(const LongTritNumber& other) {
		if (this != &other) {
			for (size_t i = 0; i < Capacity; ++i) {
				trits[i] = other.trits[i];
			}
		}
		return *this;
	}
		
	bool operator==(const LongTritNumber& other) const {
		if (trits[0] != other.trits[0]) {
			return false;
		}

		for (size_t i = 0; i < Capacity; ++i) {
			if (trits[i] != other.trits[i]) {
				return false;
			}
		}
		return true;
	}
		
	bool operator!=(const LongTritNumber& other) const {
		for (size_t i = 0; i < Capacity; ++i) {
			if (trits[i] != other.trits[i]) {
				return true;
			}
		}
		return false;
	}
		
	LongTritNumber operator+(const LongTritNumber& other) const {
		LongTritNumber result;
		Trit carry = Trit::False;

		for (int i = Capacity - 1; i >= 0; --i) {
			Trit sum = addTrits(trits[i], other.trits[i], carry);
			result.trits[i] = sum;
		}

		return result;
	}
		
	LongTritNumber operator-(const LongTritNumber& other) const {
		LongTritNumber result;
		Trit borrow = Trit::False;
		
		for (int i = Capacity - 1; i >= 0; --i) {
			Trit diff = subtractTrits(trits[i], other.trits[i], borrow);
			result.trits[i] = diff;
			borrow = calculateBorrow(trits[i], other.trits[i], borrow);
		}
		
		return result;
	}
		
	int operator<(const LongTritNumber& other) const {
		if (trits[0] != other.trits[0]) {
			return (trits[0] == Trit::True);
		}
		for (size_t i = 1; i < Capacity; ++i) {
			if (trits[i] < other.trits[i]) {
				return true;
			}
			else if (trits[i] > other.trits[i]) {
				return false;
			}
		}
		return 3;
	}

	int operator>(const LongTritNumber& other) const {
		if (trits[0] != other.trits[0]) {
			return (trits[0] == Trit::False);
		}
		for (size_t i = 1; i < Capacity; ++i) {
			if (trits[i] > other.trits[i]) {
				return true;
			}
			else if (trits[i] < other.trits[i]) {
				return false;
			}
		}
		return 3;
	}

	bool operator<=(const LongTritNumber& other) const {
		if (trits[0] != other.trits[0]) {
			return (trits[0] == Trit::True);
		}
		for (size_t i = 1; i < Capacity; ++i) {
			if (trits[i] <= other.trits[i]) {
				return true;
			}
			else if (trits[i] >= other.trits[i]) {
				return false;
			}
		}
		return false;
	}

	bool operator>=(const LongTritNumber& other) const {
		if (trits[0] != other.trits[0]) {
			return (trits[0] == Trit::False);
		}
		for (size_t i = 1; i < Capacity; ++i) {
			if (trits[i] >= other.trits[i]) {
				return true;
			}
			else if (trits[i] <= other.trits[i]) {
				return false;
			}
		}
		return false;
	}

	friend ostream& operator<<(ostream& os, const LongTritNumber& number) {
		if (number.trits[0] == Trit::True) {
			os << "Отрицательное число: ";
		}
		else {
			os << "Положительное число: ";
		}
		for (size_t i = 1; i < number.trits.size(); ++i) {
			switch (number.trits[i]) {
			case Trit::False:
				os << "False";
				break;
			case Trit::Unknown:
				os << "Unknown";
				break;
			case Trit::True:
				os << "True";
				break;
			}
			os << " ";
		}
		return os;
	}

	friend istream& operator>>(istream& is, LongTritNumber& number) {
		int sign = 0;
		cout << "Введите знак числа (0 - положительное / 1 - отрицательное): ";
		is >> sign;

		if (sign != 0 && sign != 1) {
			cout << "Некорректный ввод знака числа." << endl;
			return is;
		}

		number.trits[0] = (sign == 1) ? Trit::True : Trit::False;

		for (size_t i = 1; i < number.trits.size(); ++i) {
			int userInput;
			cout << "Введите значение трита (0 - False, 1 - Unknown, 2 - True): ";
			is >> userInput;

			if (userInput >= 0 && userInput <= 2) {
				number.trits[i] = static_cast<Trit>(userInput);
			}
			else {
				cout << "Некорректный ввод, используйте значения от 0 до 2." << endl;
				return is;
			}
		}
		return is;
	}

private:
		vector<Trit> trits;

	Trit addTrits(Trit a, Trit b, Trit& carry) const {
		int sum = static_cast<int>(a) + static_cast<int>(b) + static_cast<int>(carry);
		carry = static_cast<Trit>(sum / 3);
		return static_cast<Trit>(sum % 3);
	}

	Trit calculateCarry(Trit a, Trit b, Trit carry) const {
		int sum = static_cast<int>(a) + static_cast<int>(b) + static_cast<int>(carry);
		return static_cast<Trit>(sum / 3); // 1 если sum >= 3, иначе 0
	}

	Trit subtractTrits(Trit a, Trit b, Trit& borrow) const {
		int diff = static_cast<int>(a) - static_cast<int>(b) - static_cast<int>(borrow);
		if (diff < 0) {
			diff += 3; 
			borrow = Trit::True;
		}
		else {
			borrow = Trit::False;
		}
		return static_cast<Trit>(diff);
	}

	Trit calculateBorrow(Trit a, Trit b, Trit borrow) const {
		int diff = static_cast<int>(a) - static_cast<int>(b) - static_cast<int>(borrow);
		return (diff < 0) ? Trit::True : Trit::False;
	}
};

int main() {
	setlocale(0, "");
	LongTritNumber a, b;

	cout << "Введите числа:" << endl;
	cout << "Число a:" << endl;
	cin >> a;

	cout << "Число b:" << endl;
	cin >> b;

	LongTritNumber c = a + b;
	LongTritNumber d = a - b;

	cout << endl;
	cout << "Результат сложения: " << c << endl;
	cout << "Результат вычитания: " << d << endl;

	if (a == b) {
		cout << "a равно b" << endl;
	}
	else {
		cout << "a не равно b" << endl;
	}

	if (a != b) {
		cout << "a не равно b" << endl;
	}
	else {
		cout << "a равно b" << endl;
	}

	if (a < b) {
		if ((a < b) == 3) {
			cout << "Числа равны" << endl;

		}
		else {
			cout << "a меньше b" << endl;
		}
	}
	else {		
		cout << "a не меньше b" << endl;
	}

	if (a > b) {
		if ((a > b) == 3){
			cout << "Числа равны" << endl;
		}
		else {
			cout << "a больше b" << endl;
		}		
	}
	else {
		cout << "a не больше b" << endl;
	}

	if (a <= b) {
		cout << "a меньше или равно b" << endl;
	}
	else {
		cout << "a не меньше или равно b" << endl;
	}

	if (a >= b) {
		cout << "a больше или равно b" << endl;
	}
	else {
		cout << "a не больше или равно b" << endl;
	}

	return 0;
}
\end{lstlisting}

\newpage
\begin{center}
    \section{Исследовательская часть}
\end{center}

Далее на рисунках представлен результат работы программы.

Результат работы программы, где пользователь вводит троичное число, и программа выводит сумму двух чисел, разность и сравнивает их. Результат приведён на рис.~\ref{img:grap1}. 
\begin{figure}[h!]
    \center{\includegraphics[width=18cm]{test1.PNG}}
    \caption{Пример №1}
    \label{img:grap1}
\end{figure}

\newpage
Результат работы программы, где пользователь вводит троичное число, и программа выводит сумму двух чисел, разность и сравнивает их. Результат приведён на рис.~\ref{img:grap2}. 
\begin{figure}[h!]
    \center{\includegraphics[width=18cm]{test2.PNG}}
    \caption{Пример №2}
    \label{img:grap2}
\end{figure}

Результат работы программы, где пользователь неверно вводит троичное число. Программа выводит сумму двух чисел, разность и сравнивает их. Результат приведён на рис.~\ref{img:grap3}. 
\begin{figure}[h!]
    \center{\includegraphics[width=18cm]{test3.PNG}}
    \caption{Пример №3}
    \label{img:grap3}
\end{figure}

\newpage
Результат работы программы, где пользователь вводит два отрицательных троичных числа. Программа выводит сумму двух чисел, разность и сравнивает их. Результат приведён на рис.~\ref{img:grap4}. 
\begin{figure}[h!]
    \center{\includegraphics[width=18cm]{test4.PNG}}
    \caption{Пример №4}
    \label{img:grap4}
\end{figure}

\newpage
Результат работы программы, где пользователь неверно вводит знак троичных чисел. Программа выводит сумму двух чисел, разность и сравнивает их. Результат приведён на рис.~\ref{img:grap5}. 
\begin{figure}[h!]
    \center{\includegraphics[width=18cm]{test5.PNG}}
    \caption{Пример №5}
    \label{img:grap5}
\end{figure}

\newpage
\begin{center}
    \section*{Заключение}
    \addcontentsline{toc}{section}{Заключение}
\end{center}

Цель лабораторной работы достигнута: создана собственная структура данных.

В результате выполнения лабораторной работы были выполнены все задачи.
\begin{enumerate}
\item Описаны троичные числа.
\item Разработана структура данных для хранения целых чисел с повышением ёмкости хранения по сравнению с базовым типом данных.
\item Реализованы $6$ операторов сравнения, операторы присваивания (копии значения), сложения и вычитания двух структур длинных чисел.
\item Выполнено тестирование.
\end{enumerate}

\newpage
\begin{center}
\begin{thebibliography}{}
\addcontentsline{toc}{section}{Список используемых источников}
\bibitem{book}Рамиль Альварес Х. Алгоритмы троичной арифметики. – М.,Фонд
«Новое тысячелетие», 2012. – 20 с.
\end{thebibliography}
\end{center}

\end{document}
